\section{Cross-Case Analytical Matrix}

The point of a case library is not to tell four stories. It is to create a
reusable classification surface.

\begin{center}
\begin{tabularx}{\textwidth}{@{}l l X X@{}}
\toprule
Case & Dominant variables & Path dependence / emergence & Type \\
\midrule
Study start &
high $\Reward$, high $\Switch$, high $\Env$ &
post-lunch cue repeatedly rebuilds distraction regime &
re-entry failure \\

Exercise non-start &
low $\Will$, high barrier, route-dependent $\Env$ &
commute repeatedly selects the home-rest sequence &
boundary failure \\

Writing avoidance &
high ambiguity, moderate pseudo-work reward &
adjacent low-exposure work becomes historically efficient &
neighbouring-basin failure \\

Sleep breakdown &
high night reward, weak target-state identity, strong $\Env$ &
repeated late-night stimulation stabilizes a delayed shutdown regime &
replacement-state failure \\
\bottomrule
\end{tabularx}
\end{center}

To preserve the analytical content without overloading the table, the
cross-case dynamics and repair logic are summarized immediately below.

\begin{itemize}[nosep]
  \item \textbf{Study start.} Dynamics language: rapid re-entry into a familiar
        basin plus weak work-state restabilization. Highest-value repair:
        preload the exact next task and redesign the post-break surface.
  \item \textbf{Exercise non-start.} Dynamics language: boundary-sensitive
        trajectory selection; home acts like a basin channel. Highest-value
        repair: move the decision upstream and reroute entry before home
        capture.
  \item \textbf{Writing avoidance.} Dynamics language: movement inside a
        neighbouring attractor rather than total inactivity. Highest-value
        repair: shrink the output unit and begin with visible production.
  \item \textbf{Sleep breakdown.} Dynamics language: regime persistence and
        transition-threshold sensitivity. Highest-value repair: install a
        positive shutdown state and protect the first cue.
\end{itemize}
